\section{Reproducibility}
\label{app:repro}
Every number in Sections~\ref{sec:structural} and \ref{sec:agents} is a macro generated by \texttt{scripts/make\_paper\_assets.py} from \texttt{results/structural.json} and \texttt{results/reference\_tests.json}. The command \texttt{bash reproduce.sh} fetches DevBench at commit \texttt{\DevBenchCommit}, runs the unit tests and regenerates all results, and two runs give byte-identical output. Agent runs store every prompt and response with its hash, so the agent tables can be rebuilt from the cache without calling any model. The pre-registration file and its version-control tag are part of the released code.

\begin{table}[h]
\centering
\caption{DevBench reference code tested against its own tests (Python projects, fresh virtual environment, coverage flags removed).}
\label{tab:ref}
\small
\IfFileExists{generated/table_reference.tex}{\begin{tabular}{@{}lrr@{}}
\toprule
Project & Unit & Acceptance \\
\midrule
ArXiv\_digest & 38/38 & 0/1 \\
Hybrid\_Images & 11/19 & 1/1 \\
TextCNN & \multicolumn{2}{l}{skipped here; 5/13 on Colab} \\
chakin & 0/1 & 0/1 \\
geotext & 4/4 & 5/5 \\
hone & 7/7 & 5/5 \\
particle-swarm-optimization & 5/5 & 1/1 \\
readtime & 8/8 & 4/4 \\
stocktrends & 7/7 & 2/2 \\
lice & 25/25 & 6/6 \\
\bottomrule
\end{tabular}
}{\pending{reference tests}}
\end{table}

\section{Duty-Graph Construction Details}
Mermaid class diagrams are parsed with a tolerant parser that handles class bodies, \texttt{Class : member} lines, generic names such as \texttt{Queue\textasciitilde T\textasciitilde}, cardinalities and labels. Overloaded methods are merged into one duty. Relations to classes that are only mentioned, such as library types, are ignored because no agent implements them, and duplicate relation declarations count once.
